\subsection {W-4}
CAP teorém a jeho vztah k NoSQL databázím.

\subsubsection*{Definice a původ}

\textbf{CAP teorém}, známý také jako \textbf{Brewerův teorém}, je princip v teorii distribuovaných systémů. Tvrdí, že jakýkoli distribuovaný datový úložiště může současně zajistit pouze dvě z následujících tří vlastností:

\begin{itemize}
    \item \textbf{Konzistence (C):} Každé čtení vrací buď nejnovější záznam, nebo chybu. (Pozor: konzistence v CAP je jiná než konzistence v ACID.)
    \item \textbf{Dostupnost (A):} Každý požadavek zaslaný neporušenému uzlu musí vést k odpovědi (bez chyby).
    \item \textbf{Tolerování rozdělení (P):} Systém funguje navzdory libovolnému počtu ztracených nebo zpožděných zpráv mezi uzly v síti.
\end{itemize}

\subsubsection*{Vysvětlení principu}

Pokud dojde k \textbf{rozdělení sítě (network partition)}, je třeba se rozhodnout mezi:

\begin{itemize}
    \item \textbf{Zachování konzistence:} Operace se zruší (chyba nebo timeout), aby byla zajištěna konzistence.
    \item \textbf{Zachování dostupnosti:} Operace proběhne, ale může dojít k dočasné nekonzistenci.
\end{itemize}

\subsubsection*{Důsledek:}

Pokud dojde k výpadku sítě, nelze současně zachovat konzistenci i dostupnost. Systém musí jednu z těchto vlastností obětovat.

\subsubsection*{Normální stav (bez rozdělení)}

V případě, že k rozdělení sítě nedojde, může systém poskytovat jak konzistenci, tak dostupnost zároveň.

\subsection*{ACID vs. BASE: rozdíl mezi SQL a NoSQL databázemi}

\subsubsection*{SQL databáze (relační)}

Relační databáze (např. PostgreSQL, MySQL, Oracle) používají tradiční model \textbf{ACID}:

\begin{itemize}
    \item \textbf{A - Atomicita:} Transakce se provádí celá, nebo vůbec.
    \item \textbf{C - Konzistence:} Data musí být vždy v konzistentním stavu.
    \item \textbf{I - Izolace:} Paralelní transakce se navzájem neovlivňují.
    \item \textbf{D - Trvalost (Durability):} Po potvrzení jsou data trvale uložena.
\end{itemize}

ACID zajišťuje silnou konzistenci a je vhodný pro systémy, kde je integrita dat klíčová (např. bankovnictví, účetnictví).

\subsubsection*{NoSQL databáze}

NoSQL databáze (např. Cassandra, MongoDB, CouchDB) typicky implementují model \textbf{BASE}:

\begin{itemize}
    \item \textbf{B - Basically Available:} Systém poskytuje odpověď na každý požadavek, i když některá data mohou být momentálně nedostupná.
    \item \textbf{S - Soft state:} Stav systému se může měnit bez vstupu uživatele, kvůli asynchronní replikaci.
    \item \textbf{E - Eventual consistency:} Systém dosáhne konzistence v čase, ale ne okamžitě.
\end{itemize}

BASE obětuje okamžitou konzistenci pro vyšší dostupnost a škálovatelnost, což je vhodné pro webové aplikace, sociální sítě, analytické systémy a další rozsáhlé systémy.

\subsubsection*{Shrnutí rozdílů}

\begin{center}
\begin{tabular}{|l|l|l|}
\hline
\textbf{Vlastnost} & \textbf{SQL (ACID)} & \textbf{NoSQL (BASE)} \\
\hline
Konzistence & Silná & Postupná (eventual) \\
\hline
Dostupnost & Méně důležitá & Vysoce preferovaná \\
\hline
Transakce & Ano & Omezené nebo žádné \\
\hline
Škálovatelnost & Vertikální & Horizontální \\
\hline
Použití & Kritické systémy & Velké distribuované systémy \\
\hline
\end{tabular}
\end{center}

\subsection*{Vztah k CAP teorému}

\begin{itemize}
    \item \textbf{SQL databáze (ACID)}: upřednostňují \textbf{Konzistenci} před Dostupností $\Rightarrow$ typicky \textbf{CP}.
    \item \textbf{NoSQL databáze (BASE)}: upřednostňují \textbf{Dostupnost} před Konzistencí $\Rightarrow$ typicky \textbf{AP}.
\end{itemize}

\subsection*{Odpovědi na otázky}

\subsubsection*{Proč není možné mít v distribuovaném systému všechny tři vlastnosti CAP současně?}

Protože během výpadku komunikace (partition) musí systém buď:
\begin{itemize}
    \item odmítnout operace (ztrácí dostupnost), nebo
    \item provést operace bez záruky aktuálnosti dat (ztrácí konzistenci).
\end{itemize}

\subsubsection*{Jaký dopad má CAP teorém na výběr databáze pro kritický finanční systém?}

Finanční systémy vyžadují silnou konzistenci. Je proto vhodné použít \textbf{ACID} databázi, která implementuje model \textbf{CP} (např. PostgreSQL, Oracle, HBase).

\subsubsection*{Vysvětlete rozdíl mezi eventual consistency a strong consistency.}

\begin{itemize}
    \item \textbf{Strong consistency:} Všichni klienti vždy čtou nejnovější záznam.
    \item \textbf{Eventual consistency:} Klienti mohou číst starší hodnoty, ale data se časem sjednotí.
\end{itemize}

\subsubsection*{Proč je partition tolerance u moderních NoSQL databází považována za nevyhnutelnou?}

Protože v distribuovaném prostředí jsou výpadky a ztráty zpráv nevyhnutelné. Aby systém zůstal funkční, musí být schopen tyto poruchy tolerovat $\Rightarrow$ \textbf{Partition Tolerance (P)} je nutností.
